% 07_budget.tex - Presupuesto de la propuesta.
\section{Presupuesto}
% SSOT del presupuesto: data/budget.yaml (editado por research-architect con
% los importes que de el investigador). REGLA DURA: ninguna cifra economica
% en esta seccion procede de otra fuente; mientras el investigador no la
% aporte, TODA cifra es \todo{DATO PENDIENTE}. JAMAS inventar euros.

\subsection{Presupuesto por categorias}
% TODO: tabla derivada de data/budget.yaml. Importes pendientes hasta que el
% investigador los aporte; el desglose por socio/categoria debe cuadrar con
% el total del YAML.
\begin{table}[htbp]
\centering
\caption{Presupuesto por categorias (derivado de data/budget.yaml).}
\label{tab:budget}
\begin{tabular}{lll}
\toprule
Categoria & Descripcion & Importe (EUR) \\
\midrule
Personal & TODO: personal contratado y dedicaciones & \todo{DATO PENDIENTE: personal} \\
Costes externos & TODO: subcontrataciones, servicios y viajes & \todo{DATO PENDIENTE: costes externos} \\
Equipamiento & TODO: equipos e infraestructura & \todo{DATO PENDIENTE: equipamiento} \\
Costes indirectos & TODO: segun las bases de la convocatoria & \todo{DATO PENDIENTE: costes indirectos} \\
\midrule
Total & & \todo{DATO PENDIENTE: presupuesto total} \\
\bottomrule
\end{tabular}
\end{table}

\subsection{Justificacion del presupuesto}
% TODO: justificacion de cada categoria (por que es necesaria para los WPs y
% coherente con data/workpackages.yaml). Reglas de elegibilidad y maximos de
% la convocatoria: context/call_guidelines.md y
% context/convocatoria_checklist.md.

\subsection{Aportacion propia del socio privado}
% TODO: si las bases lo piden, cofinanciacion o aportacion en especie del
% socio privado (cuantia pendiente de decision del investigador:
% \todo{DATO PENDIENTE: cofinanciacion}).
